Nog immer was Emmy
allerminst een lesbie.
In de heks haar tronie,

hoezeer ook haar afkeer,
zag ze dat zij weleer
denkbaar oogde tot meer.

Met rimpels op leeftijd,
de spataderen wijd
en tumoren verspreid,

was de schoonheid vergaan.
Haar laatste eigenwaan
was door de tijd ontdaan.

De zwaarste verminking
kwam door de bezinking
van Jeanne’s verdrinking

in valse intriges
geleid door prestiges.
Haar zwakke voltiges

gaf vijanden het zuur.
Het was van korte duur,
maar sloeg een vitaal uur.

De littekens bleven.
Ze kon niet vergeven
dat haar mond bleef kleven.

De grijze haren op
deze Michelin pop
liepen vanaf haar kop

tot over Jeanne’s pruim.
Haar gleuf was uiterst ruim
en droop steeds van het schuim

dat soms naar buiten spoot.
Emmy was in watersnood
als haar mond het omsloot.

Gore verrichtingen,
strenge verplichtingen
in vele richtingen

werden van haar verwacht
en tot uitvoer gebracht
zonder enige klacht.

Emmy hield schoon het huis alsmede Jeanne’s muis.
Emmy droeg braaf haar kruis

als de heks doordramde.
De heks ontvlamde
als Emmy haar kamde

zonder borstelwaren.
Alle grijze haren
zelfs de onzichtbaren,

bijvoorbeeld in haar neus,
streek Emmy minutieus.
Dat bleek niet delicieus

want het moest met de tong. De heks aan het snotjong
aldus haar drek opdrong.

Denk aan neusslijm of talg.
“De heks weet dat ik walg.
Ze verdient dus de galg,

maar wordt nooit gehangen”.
Emmy bleef gevangen
in de vele gangen

van schilfers, zweet, oorsmeer,
eelt of pus uit een zweer.
Eerder geslachtsverkeer

passeerde de revue
als gang op het menu.
De heks hield continu

zonder genade thuis.
Emmy at per abuis
zelfs een keer een schaamluis.

Als de heks op haar beurt
Emmy, door haar besmeurd,
somtijds goedgehumeurd

plots wilde verrassen,
met háár muis te wassen
door erop te plassen,

dan ging Emmy tobben.
Want het werd gauw schrobben
om Emmy te mobben.

De knokige washand
werd dan een penetrant,
die het pis langs de rand

als glijmiddel aangreep.
Met urine als zeep
zich aan Emmy vergreep.

Zo bracht de huisbazin
haar hele vuist erin.
“Ben ik je koningin?",

polste de dwingeland
met haar dwingende hand.
Emmy knikte galant

en kuste toen de vuist,
gebald en nat, die juist
uit haar muis was verhuisd.

Emmy moest naast zoenen
ook zelf stevig boenen.
Er waren miljoenen

te schoonmaken spullen.
De heks had veel prullen
om haar huis te vullen.

Het huis schoonmaken
moest ook veroorzaken
haar seksuele taken.

Want als Emmy geknield
slaafs een vloer onderhield
werd ze weer snel vernield.

Haar wiebelende kont
ontbloot boven de grond
ontstak de heks haar lont.

Die aanblik was bedoeld;
“dat de heks weer wat voelt
na te zijn afgekoeld”

zo ontdekte Emmy.
Dat was de tragedie
voor de rode sloerie.

De vervreemde Emmy
hoorde over magie
tijdens de liturgie,

maar leerde wat het was
in haar eenpersoons klas
op de heks haar matras.

Want de heks maakte haar
door misbruik gebruiksklaar
als wufte handelswaar

voor haar tovenarij.
Zo stond ze naast haar zij
al roerend door de brij

van bloed en brandnetel
in een heksenketel
en gaf haar een zetel.

Een kikker nog levend
moest Emmy al bevend,
droevig en meelevend

in het brouwsel koken,
door het vuur te poken
en meer hout te stoken.

Was de soep eenmaal gaar
dan duwde de heks haar
gezicht in het etenswaar.

Emmy stikte zowat
tot ze de kikker at
glijdend door haar keelgat.

De vegetariër
werd door de senior
vervormd tot magiër.

Iedere weifeling,
elke vertwijfeling
leidde tot bestraffing.

Dat was precies het plan.
De heks kneep Emmy dan
en zo viel in de pan

alsnog een levend dier.
Dit vleesminnend geklier
schonk de heks veel plezier.

Niet ontmenselijking
was waar ze met haar “ding” mee over de schreef ging,

maar bij een stukke lip
of een gebroken rib
van zijn tere broedkip

kon ze heibel voorzien.
Het bestond bovendien
dat Emmy dan misschien

haar mondje niet meer hield.
Of zo erg was vernield
dat ze hem seks onthield.

Aan de jager melden
dat minnaars haar kwelden
zou Emmy ontgelden.

Eens bij haar kutscheten,
die Emmy moest “eten”,
liet de heks haar weten

dat zij ooit had gemoord,
toen zij niet werd gehoord.
De vrucht van dit machtswoord

spoorde haar nog meer aan
om er op in te gaan,
hoe ze het had gedaan.

Het was in Emmy's kerk.
Het was stevig handwerk,
want de luchtpijp was sterk.

De verzwakte Heidi
moest dood als retorsie
voor haar depositie

over haar handlanger.
Emmy, toen al zwanger,
was een dromenvanger

toen wraak werd genomen.
Ze had het vernomen
in haar wildste dromen.

De priester als trawant
had een duistere kant,
die hij met harde hand

in Emmy uitleefde
en ze vaak beleefde.
Waarvan Emmy beefde

was dat een kerkganger
zowaar een aanhanger
bleek van een katvanger,

die de duivel diende.
Jeanne hielp haar vrienden
als ze het verdienden.

De priester bekwam het.
Hij zorgde ooit maar net
dat ze werd uitgezet.

De buur die ze veracht,
kreeg aldus geen volmacht
op een pikante nacht.

Deze openbaring
plus haar gewaarwording
bracht de overtuiging

dat de pil te slikken
om te gaan verklikken
haar zou doen verstikken.

Daarnaast was er de kwaal.
De erfelijke taal,
die door Jeanne’s gemaal

niet kan worden ontkend.
Hoe ze hem ook verwend,
haar kind lijkt op de vent,

die de jager zo haat.
Ze zag als juiste daad
te dansen met het kwaad.

Door Emmy’s overspel,
daartoe genoopt dat wel,
zat ze flink in de knel.

Emmy bleek al vaker
bij een hersenkraker
een brommende braker.

De migraineaanval
duurde meestal nogal.
Verschillend per geval.

Deze onderbreking
was een ontgoocheling
wat betreft de planning.

Dus gebruikte Jeanne
de grond van verbannen
voor haar snode plannen.

Ze gaf tegen hoofdpijn
en kotsmisselijk zijn,
gember als medicijn.

Gingerol bleek funest
voor Emmy’s stil protest
tegen een valse test.

Zo kon Jeanne even
zonder medeleven
zich op haar uitleven.
